% 第5章 5.4.1：Mercer 特征函数与特征值衰减。
% (a) 最小值核 min(x,z) 在 [0,1]、Lebesgue 测度下的前三个特征函数 psi_j=sqrt2 sin((j-1/2)pi x)（解析式）。
% (b) 特征值（对数纵轴）：最小值核 1/((j-1/2)^2 pi^2)（解析）；高斯核 sigma=0.5、[-1,1]、dz/2
%     （同目录 generate.py 用 40 点 Gauss-Legendre 离散，与 60 点结果前12位一致到 1e-3 相对精度）；仿射核 1、1/3。
\begin{tikzpicture}[font=\figurefont\footnotesize]
\pgfplotsset{every axis/.style={width=52mm,height=40mm,axis lines=left,
  x axis line style={draw=FigureStroke,line width=.6pt},y axis line style={draw=FigureStroke,line width=.6pt},
  tick style={draw=FigureStroke},tick label style={font=\figurefont\scriptsize,text=FigureInk},
  label style={font=\figurefont\scriptsize},
  ymajorgrids,grid style={draw=FigureGrid,line width=.4pt},
  title style={font=\figurefont\footnotesize},
  legend style={draw=none,fill=white,fill opacity=.85,text opacity=1,font=\figurefont\scriptsize,cells={anchor=west}}}}
\begin{axis}[name=A,xmin=0,xmax=1,ymin=-1.6,ymax=1.6,xtick={0,0.5,1},ytick={-1.414,0,1.414},yticklabels={$-\sqrt2$,$0$,$\sqrt2$},
  xlabel={$x$},xlabel style={at={(axis description cs:1,0.5)},anchor=west},
  title={(a) 最小值核的特征函数},samples=121,domain=0:1,
  legend style={at={(0.03,0.03)},anchor=south west}]
\addplot[draw=FigureBlue,line width=1.1pt] {sqrt(2)*sin(deg(0.5*pi*x))};
\addplot[draw=FigureGreen,line width=1.1pt,dashed] {sqrt(2)*sin(deg(1.5*pi*x))};
\addplot[draw=FigureYellow,line width=1.1pt,dash dot] {sqrt(2)*sin(deg(2.5*pi*x))};
\draw[FigureStroke,line width=.5pt] (axis cs:0,0) -- (axis cs:1,0);
\legend{$\psi_1$,$\psi_2$,$\psi_3$}
\end{axis}
\begin{semilogyaxis}[at={(A.east)},anchor=west,xshift=11mm,xmin=0.5,xmax=12.5,ymin=1e-9,ymax=2,
  xtick={1,3,5,7,9,11},ytick={1,1e-2,1e-4,1e-6,1e-8},xlabel={$j$},
  xlabel style={at={(axis description cs:1,0)},anchor=north west,yshift=2mm},
  title={(b) 特征值$\lambda_j$},
  legend style={at={(0.03,0.03)},anchor=south west}]
\addplot[FigureBlue,line width=.8pt,mark=*,mark size=1.3pt] coordinates {(1,4.052847e-01) (2,4.503164e-02) (3,1.621139e-02) (4,8.271117e-03) (5,5.003515e-03) (6,3.349461e-03) (7,2.398135e-03) (8,1.801265e-03) (9,1.402369e-03) (10,1.122672e-03) (11,9.190130e-04) (12,7.661337e-04)};
\addplot[FigureGreen,line width=.8pt,dashed,mark=square*,mark size=1.2pt,mark options={solid}] coordinates {(1,5.212157e-01) (2,3.025390e-01) (3,1.257423e-01) (4,3.892388e-02) (5,9.388322e-03) (6,1.839349e-03) (7,3.025908e-04) (8,4.285715e-05) (9,5.325958e-06) (10,5.893754e-07) (11,5.876784e-08) (12,5.331406e-09)};
\addplot[only marks,FigureYellow,mark=triangle*,mark size=2pt] coordinates {(1,1) (2,0.333333)};
\legend{最小值核,高斯核,仿射核}
\end{semilogyaxis}
\end{tikzpicture}
